\documentclass[11pt]{article}
\usepackage{irvlreport}

% --- packages the body needs -----------------------------------------------
\usepackage{amsmath}      % equation, \text{...} inside math
\usepackage{booktabs}     % \toprule \midrule \bottomrule
\usepackage{graphicx}     % \includegraphics
\usepackage{url}          % \url{} inside the .bib entries
\usepackage{listings}     % \lstinputlisting for the appendix
\usepackage{natbib}       % \citep and \citet

\graphicspath{{figures/}{./}}
\bibliographystyle{plainnat}

\lstset{basicstyle=\ttfamily\footnotesize,breaklines=true,columns=fullflexible}

% --- local macro used throughout the draft ---------------------------------
% Marks a value still to be verified. Delete the definition and every call
% before submission.
\newcommand{\placeholder}[1]{\textbf{[#1]}}

\irvlsettitle{From Reasoning to Action: Benchmarking Frontier AI Models on LIBERO-10 Robot Manipulation Tasks}
\irvlsetshorttitle{LIBERO-10 Benchmark}
\irvlsetauthors{May Espinola\textsuperscript{*}, Qifan Zhang\textsuperscript{*}, Yu Xiang\\[3pt]
\normalsize\textsuperscript{*}Equal contribution}
\irvlsetdate{October 2026}
\irvlsetversion{Version 1.0}
\irvlsetlogo{irvl-logo.png}
% Optional: \irvlsetreportid{IRVL-TR-2026-01}
\irvlsetprojecturl{https://irvlutd.github.io/LLM-LIBERO/}
\irvlsetcodeurl{https://github.com/IRVLUTD/LLM_LIBERO}

\begin{document}
\irvlmaketitle
\begin{irvlabstract}
Multimodal language models can interpret visual observations and use tool calls, which makes it possible to evaluate them as robot controllers without training a task-specific policy. Such an agent also declares its own completion, so an evaluation must score the resulting environment state rather than the agent's report. We study this setting on the ten simulated manipulation tasks of LIBERO-10 under a single manipulation framework shared by every API-served model we test. Each agent receives a natural-language goal, two RGB camera views, and robot proprioception, and controls a Franka Panda arm through absolute end-effector pose commands. The protocol holds the observation interface, manipulation prompt, simulator horizon, and model-call allowance fixed while varying the model and reasoning effort. We derive the per-episode allowance of 80 model calls from the call usage of successful pilot trajectories with a factor-of-two margin. There are fifteen configurations total: GPT-6.1 Sol, GPT-6 Astra, GPT-6 Luna, Claude Opus 5, and Claude Opus 5.5, crossed with three requested reasoning levels (low, medium, high). For every episode we record LIBERO's own success predicate, the termination category, model calls and simulator steps consumed, input, cached, output and reasoning tokens, API latency, and estimated cost, both per call and in aggregate. Per-episode records and rollout videos are browsable at the project site.
\end{irvlabstract}

% ===========================================================================
\section{Introduction}
% ===========================================================================
Language-guided manipulation requires identifying objects, planning motions,
and checking whether an action produced the intended physical outcome. An
API-served multimodal language model can attempt all three by alternating
visual observations with robot-control tool calls. Producing a plausible plan
or declaring completion, however, establishes nothing about the resulting
physical state. An evaluation of these models as controllers must therefore
score the environment rather than the agent's report, and must account for the
inference resources the attempt consumed.

We investigate how model identity and requested reasoning effort affect
manipulation under a common closed-loop interface. The protocol has three
fixed points: every configuration sees one observation and action interface,
success is the simulator's own predicate rather than the agent's report, and
each episode runs under a declared allowance of model calls whose cost is
accounted for and reported. We use LIBERO-10, drawn from the LIBERO
benchmark~\citep{liu2023libero}, and connect the simulator to the upstream
Inspect Robots agent~\citep{inspectrobots} through a custom embodiment
adapter. The agent selects pose targets directly; no task-specific policy is
trained. Our study concerns zero-shot execution under general robot
instructions, not lifelong policy learning.

Three recent studies evaluate this same generation of API-served models as
embodied controllers, and together they locate our design.
\citet{galbot2026astra} evaluate GPT-6 Astra across six embodiment domains and
report a gap between competent task-level decisions and reliable physical
control: in their \emph{Direct} configuration, where the model emits control
targets itself, Astra reaches 26\% on the 50-instance RoboDojo split against
48\% when it instead reviews and corrects a learned policy's proposals.
\citet{roboicl2026} hold that model fixed and change only its context, raising
the RoboDojo progress score from 33.68 to 50.64 by writing demonstrations in
the same observation--action--receipt--observation grammar as the agent's own
interaction history. \citet{vlmagents2026} likewise keep the model fixed and
vary the family of demonstration supplied through a context compiler.

Our protocol is deliberately the \emph{Direct}, zero-shot, single-interface
corner of that space: one frozen model, no learned policy in the loop, no
demonstrations, and one observation and action interface shared across every
configuration. This is the weakest configuration surveyed above, and we choose
it for two reasons. It isolates what the model contributes on its own, and it
establishes the no-demonstration reference point that context-engineering
results are measured against but, on LIBERO, do not yet report. Each of those
studies also evaluates a single model, so holding the interface constant across
models is what makes the complementary comparison possible here.

The study asks three questions: whether these models can complete multi-stage
manipulation from RGB and proprioception; how reasoning effort affects success
and inference cost; and which observable failures limit execution. The
contribution of this report is a documented protocol with its design rationale,
and an auditable record of observations, tool calls, simulator actions, and API
usage.

% ===========================================================================
\section{LLM-based Manipulation Framework}
\label{sec:method}
% ===========================================================================
At decision step $k$ the model receives the task instruction, the current
observation, and the retained conversation history, and returns exactly one
tool call. For a motion call, the framework interpolates the requested target
into low-level actions, executes them, and returns a new observation. One model
call therefore advances many simulator control steps, and the model must
inspect the next observation to judge whether the target was reached and
whether a grasp or placement succeeded. Figure~\ref{fig:interface} shows one
such step.

\begin{figure}[t]
\centering
% \graphicspath already points at figures/ -- a path prefix here will break
% the build from any other directory.
\includegraphics[width=\textwidth]{LIBERO_Framework.pdf}
\caption{\textbf{Overview of the LLM-based manipulation framework in LIBERO.}
A multimodal LLM uses previous and current images, robot state, and prompt context to select \texttt{move\_to}, \texttt{done}, or \texttt{give\_up}.
The adapter converts \texttt{move\_to} targets into control actions, and updated observations inform the next decision.
Interaction ends upon \texttt{done}, \texttt{give\_up}, or a call/step limit.
LIBERO's success checker independently determines task success.}
\label{fig:interface}
\end{figure}

\subsection{Observations}
Each observation contains an external agentview image and a wrist-camera image,
both $256\times256$ RGB. Proprioception comprises world-frame end-effector
position, a six-dimensional orientation representation, the current gripper
command, and measured finger joint positions.

Four choices here are deliberate, and each would be a confound if left
unstated. First, we supply two views rather than one: the agentview camera
alone leaves the gripper's relation to a target ambiguous at contact, which is
where LIBERO-10's multi-stage tasks are decided, and the wrist view is what
makes a pre-grasp alignment checkable from the observation itself. Second,
orientation is the first two columns of a rotation matrix rather than a
quaternion or Euler angles. This representation is continuous and free of the
sign ambiguity and convention choices that make quaternions and Euler triples
easy to emit in an inconsistent frame, and the third column is recoverable by
cross product. Third, we report the gripper \emph{command} and the measured
finger joint positions separately, and the prompt states that they differ,
because a closed command with the fingers still at their commanded width is the
signature of a failed grasp --- the single most consequential distinction the
agent has to make, and one the command alone cannot express. Fourth, the model
receives no privileged object coordinates, so observed pose errors are
perception-and-reasoning errors under this interface rather than arithmetic on
ground truth.

Image context retains the two most recent observations, at most four camera
images, while earlier text and tool history remain available. This limit bounds
per-request image cost, which otherwise grows linearly in retained frames
across an 80-call episode, and it fixes the visual history every configuration
sees. It also makes the interface memoryless in images beyond one step, so a
task whose later stages depend on a visual fact established early must be
solved through the retained text history. Comparable recent work sits at the
other end of that tradeoff, with image budgets near 48 per
request~\citep{roboicl2026,vlmagents2026}; we state our limit rather than treat
it as an incidental setting.

\subsection{Tools and motion execution}
The agent has three tools: \texttt{move\_to}, \texttt{done}, and
\texttt{give\_up}. A \texttt{move\_to} call specifies an absolute world-frame
position, optional orientation entries, and an optional gripper command;
omitted dimensions retain their current values. Each motion carries a short
\texttt{note} describing the current observation and the reason for the action.
Notes are recorded for inspection and are never a success criterion.

The action vector is $(x,y,z,r_{00},r_{10},r_{20},r_{01},r_{11},r_{21},g)$.
Translation targets are bounded by $x,y\in[-0.6,0.6]$\,m and
$z\in[0.42,1.6]$\,m; orientation entries lie in $[-1,1]$ and $g\in[0,1]$. The
framework checks bounds and per-step action changes. Per-step limits are
0.01\,m per translation coordinate, 0.05 per orientation entry, and 0.1 for the
gripper command. The adapter orthonormalizes the orientation columns before
converting to the simulator representation. The gripper is an open/close
actuator and does not accept a partial opening. Low-level control runs at
20\,Hz.

We expose \emph{absolute} pose targets with framework-side interpolation rather
than per-step deltas for two reasons. An absolute target is independently
checkable: the agent can compare the commanded pose against the proprioception
it receives next and detect a shortfall, which a sequence of deltas does not
afford unless the agent integrates its own history. And interpolation moves
responsibility for smoothness and rate limiting off the model, so a failure is
attributable to target selection rather than to an inability to emit a
well-conditioned trajectory at 20\,Hz. The cost is that one model call advances
many control steps, which makes the call allowance and the simulator horizon
separate resources that must both be reported; \citet{galbot2026astra} adopt
the same separation, distinguishing control steps, action segments, and model
requests. The per-step limits are the simulator-side rate limits that make
interpolation physical, bounding how far a single accepted target can move the
arm per control step.

\subsection{Prompt design}
Every configuration uses one prompt structure. The system message describes the
tool-control loop, requires exactly one tool call per turn, requests an
explanatory note for each motion, and states the 80-call allowance. Embodiment
notes describe coordinates, orientation, gripper semantics, the available
observations, and general manipulation guidance: maintain clearance,
distinguish distractors, verify a grasp with a short lift, check the result of
each motion, and change strategy after repeated failures. The user message
supplies the native task instruction as \texttt{Goal: <instruction>}.

The zero-shot condition supplies no demonstrations, no demonstration-derived
task notes, and no learnings carried from earlier episodes. Each episode starts
a fresh policy instance. The terminal tools request a retrospective
observation, which is logged but never passed to a later episode. Zero-shot
here means execution without task demonstrations or task-specific training, not
execution without general manipulation guidance; the guidance is part of the
interface and is held constant, so it cannot explain a difference between
configurations. Appendix~\ref{app:prompt} reproduces the archived system
prompt.

% ===========================================================================
\section{Experimental Setup}
\label{sec:experiments}
% ===========================================================================

\subsection{Tasks and initialization}
The evaluation covers all ten native LIBERO-10 tasks, indexed 0--9 and listed
in Appendix~\ref{app:tasks}; they include multi-object placement, stove
interaction, drawer closure, and microwave placement and closure. Each task
provides 50 saved initial states, and the full batches use states 0--9, giving
100 episodes per model--effort configuration.

The ten-states-per-task design trades repetitions for initial-state coverage.
Variation across initial states is a property of the benchmark that we want
represented in the estimate; variation across repeated calls on the same state
is provider sampling noise. Spending a task's episode allowance on ten distinct
states rather than on repetitions of fewer states therefore targets the
variance component that generalizes. The cost is that each (task, state) cell
is observed once per configuration, so cell-level outcomes carry no
within-cell variance estimate, and we quote uncertainty only at the task and
configuration level.

For state index $i$, the simulator loads the native saved state with
environment seed $142+i$. Ten neutral steps with an open-gripper command settle
the scene before the scored step counter is reset. The evaluation-framework
seed is $42+e$, where $e$ is the episode index within a task batch.
Simulator seeds do not make remote model generation deterministic, so repeated
episodes on one state are not expected to be identical. Runs proceed
sequentially in configuration blocks rather than in randomized or
counterbalanced order, so a provider-side change over the campaign period would
not be distributed evenly across configurations; run timestamps are archived so
that this can be checked.

\subsection{Configuration coverage}
\label{sec:coverage}
The campaign is a full factorial of five models crossed with three requested
reasoning levels, giving the fifteen configurations of
Table~\ref{tab:coverage}. Each runs the same ten tasks from the same ten saved
initial states, so every configuration contributes 100 episodes and every
comparison is matched on task and initial state. All fifteen are complete.

\begin{table}[h]
\centering
\caption{Configuration matrix and coverage. $T$ is the number of LIBERO-10
tasks evaluated and $S$ the number of saved initial states per task; episodes
$=T\times S$. Effort labels are defined per model family in
Section~\ref{sec:models}.}
\label{tab:coverage}
\small
\begin{tabular}{llccrl}
\toprule
Model & Effort & $T$ & $S$ & Episodes & Status \\
\midrule
GPT-6.1 Sol     & low    & 10 & 10 & 100 & complete \\
GPT-6.1 Sol     & medium & 10 & 10 & 100 & complete \\
GPT-6.1 Sol     & high   & 10 & 10 & 100 & complete \\
GPT-6 Astra     & low    & 10 & 10 & 100 & complete \\
GPT-6 Astra     & medium & 10 & 10 & 100 & complete \\
GPT-6 Astra     & high   & 10 & 10 & 100 & complete \\
GPT-6 Luna      & low    & 10 & 10 & 100 & complete \\
GPT-6 Luna      & medium & 10 & 10 & 100 & complete \\
GPT-6 Luna      & high   & 10 & 10 & 100 & complete \\
\midrule
Claude Opus 5   & low    & 10 & 10 & 100 & complete \\
Claude Opus 5   & medium & 10 & 10 & 100 & complete \\
Claude Opus 5   & high   & 10 & 10 & 100 & complete \\
Claude Opus 5.5 & low    & 10 & 10 & 100 & complete \\
Claude Opus 5.5 & medium & 10 & 10 & 100 & complete \\
Claude Opus 5.5 & high   & 10 & 10 & 100 & complete \\
\bottomrule
\end{tabular}
\end{table}

All fifteen comparisons draw on the full $10\times10$ grid, so model identity
and effort are crossed rather than confounded. Per-episode
records, rollout videos, and the live state of the matrix are published at the
project site.

\subsection{Models and inference settings}
\label{sec:models}
The campaign covers GPT-6.1 Sol, GPT-6 Astra, and GPT-6 Luna --- API
identifiers \texttt{openai/gpt-6.1-sol}, \texttt{openai/gpt-6-astra}, and
\texttt{openai/gpt-6-luna} --- together with Claude Opus 5 and Claude Opus 5.5
\placeholder{(API identifiers)}. Model names and effort settings are those
recorded in the archived requests, so a provider-side change to a model behind
a stable identifier would appear as unexplained variation across run
timestamps rather than as a change of configuration.

We vary model identity and effort and hold everything else fixed because these
are the two factors a practitioner can change without re-engineering the robot
stack, and because the studies closest to
ours~\citep{galbot2026astra,roboicl2026,vlmagents2026} each hold the model
fixed and vary the context instead. The complementary design --- one interface,
several models --- is what lets us ask whether an observed limitation belongs
to a particular model or to the interface we have given all of them.

\paragraph{Effort is not a cross-provider scalar.}
Within the GPT-6 family, \texttt{low}, \texttt{medium}, and \texttt{high} are
values of a single provider-exposed effort parameter, so the ordering is
meaningful within a model and the setting is identical across models of that
family. This does not extend across providers. The Claude arms expose
deliberation as a \emph{thinking-token budget} rather than an ordinal level, so
\texttt{low}, \texttt{medium}, and \texttt{high} there denote
\placeholder{thinking budgets of $B_{\text{low}}$, $B_{\text{medium}}$, and
$B_{\text{high}}$ tokens}, chosen by us rather than named by the provider. The
same budgets are used for both Claude generations, so effort is comparable
between them but not between the Claude and GPT-6 families. We therefore treat
effort as a within-family factor, and we do not read an effort difference
across families as a difference in the same underlying variable.

\paragraph{Request settings.}
Requests use the Responses interface with \texttt{service\_tier=default} and
\texttt{store=false}. Temperature is not set explicitly. The output ceiling is
\placeholder{3{,}072} tokens per response including reasoning tokens, set above
the longest response seen in pilot inspection while remaining low enough that
one pathological generation cannot consume an episode's allowance;
truncations are recorded separately so the ceiling's effect stays visible
rather than being absorbed into the failure count. Prompt wording, camera
resolution, image history, tools, and control limits are common to all
configurations. Historical prompt variants are identified by their saved
prompt hashes.

\subsection{Choosing the model-call allowance}
\label{sec:callbudget}
The allowance must satisfy two competing requirements. It must not truncate
episodes that would otherwise succeed, or the measured success rate becomes a
function of the allowance rather than of the policy. And it must bound
per-episode inference, since cost and latency scale with calls and the campaign
is large. We therefore set the allowance from the observed call usage of
\emph{successful} trajectories and applied a fixed margin of two.

Let $\mathcal{S}_{\mathrm{pilot}}$ be the set of successful pilot episodes and
$C_j$ the number of logical model calls episode $j$ consumed through
termination. The design statistic is the mean call usage of successful
trajectories,
\begin{equation}
\bar C_{\mathrm{success}}
=\frac{1}{|\mathcal{S}_{\mathrm{pilot}}|}
\sum_{j\in\mathcal{S}_{\mathrm{pilot}}}C_j ,
\label{eq:cbar}
\end{equation}
and the allowance is
\begin{equation}
C_{\max}=\big\lceil\,2\,\bar C_{\mathrm{success}}\,\big\rceil_{10}=80\ \text{calls},
\label{eq:budget}
\end{equation}
rounded to the nearest ten. Doubling is a deliberate margin rather than a tuned
threshold. A successful trajectory spends its calls on a nominal path through
the task; an episode that recovers from a dropped grasp, re-approaches a
misaligned target, or disambiguates a distractor must re-traverse part of that
path, so the recoverable cases we most want to leave room for are exactly those
whose usage lies above the successful mean. A factor of two admits one full
additional traversal. We make no claim that 80 is optimal. It is an operating
point fixed before the campaign and held constant across every configuration,
which is the property that matters for comparability. Fixing an explicit
per-episode decision budget and reporting it as a protocol parameter is
standard in this line of work; \citet{galbot2026astra}, for instance, allow a
120-decision budget for in-hand manipulation and 250 model calls for a
30-second locomotion run.

\paragraph{Ex-post validation.}
Because the allowance was fixed in advance, its adequacy is testable against
the completed runs. Over 1,500 episodes containing 285 environment-verified
successes, successful episodes used a mean of 45.5 calls (median 44, 95th
percentile 69), and the maximum over all 285 was 79. No successful episode
reached the allowance. Ninety-nine episodes did reach it, and every one of them
was a failure. The cap therefore truncated no trajectory that the environment
predicate would have scored as a success, and the realized margin,
$C_{\max}/\bar C_{\mathrm{success}}=80/45.5\approx1.76$, is close to the factor
of two the design assumed.

\paragraph{Where the margin is thin.}
The margin is not uniform across models. Successful GPT-6 episodes sit well
inside it: Astra averages 42.2 calls (95th percentile 62, maximum 74) and Sol
43.4 (65, 69). The Claude arms spend more per success --- Opus~5 averages 47.4
calls (70, 77) and Opus~5.5 averages 53.1 (77, 79) --- so for Opus~5.5 the
allowance is close to binding, and applying Eq.~\eqref{eq:budget} to that model
alone would have given a cap near 110 rather than 80. This is a limitation of
one shared allowance set from a pilot that did not include every model, and the
Opus~5.5 success rates should be read as lower bounds.

\paragraph{What remains to be documented.}
\placeholder{The pilot sample size $|\mathcal{S}_{\mathrm{pilot}}|$, the pilot
value of $\bar C_{\mathrm{success}}$, and the pilot call-count spread must be
inserted here from the preliminary logs.} Eq.~\eqref{eq:cbar} must be computed
from the preliminary runs that informed the decision, not from the benchmark
outcomes above; the ex-post figures validate a prior choice and would be
circular if presented as its justification.

\subsection{Episode limits and financial accounting}
Each episode permits at most 80 logical model calls and 900 scored simulator
steps. At 20\,Hz the latter is 45\,s of nominal simulated time, excluding
settling and API waiting. The two are separate limits, since one tool motion
can consume several control steps, and the output-token ceiling supplies a
third inference constraint.

The limits bind different resources and follow different reasoning. The call
allowance bounds inference, as derived in Section~\ref{sec:callbudget}. The
900-step horizon bounds simulated execution and is set so that it is not the
binding constraint on a nominally executed episode: at the 0.01\,m per-step
translation limit, 900 steps permit roughly 9\,m of commanded end-effector
travel, an order of magnitude more than these tasks require, so an episode that
exhausts the horizon has been executing unproductive motion rather than running
out of room. As in \citet{galbot2026astra}, simulation is paused during
inference, so the horizon measures robot motion and excludes decision latency;
latency is reported separately.

Financial allowances are spending backstops, not an experimental factor. The
call allowance is the controlled inference resource because it is the quantity
comparable across providers: a matched dollar cap would grant a cheaper model
more decisions than an expensive one and convert a price difference into an
apparent capability difference. \placeholder{State the final per-episode dollar
allowance used for each model family.} The transport reserves estimated input
cost and the full output allowance before dispatch, then settles against
returned usage. Costs are estimated from archived token rates and cache usage
rather than reconciled against provider invoices. The campaign is consequently
\emph{not} an equal-dollar-budget comparison, and per-episode cost differences
reflect provider token rates and realized call counts rather than an imposed
budget.

\subsection{Success, termination, and recovery}
Success is LIBERO's native predicate evaluated at the end of the episode.
Calling \texttt{done}, or receiving a successful framework execution status,
does not count as task success. For task $t$ and initial state $i$, let
$s_{t,i}\in\{0,1\}$ denote the predicate. With $K$ evaluated states per task,
the balanced task-averaged success rate is
\begin{equation}
\mathrm{SR}=\frac{1}{10}\sum_{t=0}^{9}\frac{1}{K}\sum_{i=0}^{K-1}s_{t,i}.
\label{eq:sr}
\end{equation}
Because every configuration covers the same $10\times10$ grid, this balanced
rate equals the unweighted mean over its 100 episodes.

Episode records distinguish successful completion, unsuccessful termination,
simulator-horizon exhaustion, call-limit exhaustion, financial stops, API
errors, and interruptions. Some legacy labels merge the call and financial
limits into one category; where they do, we recover the distinction from the
recorded call count and settled cost. Missing and interrupted outcomes are
shown explicitly rather than silently scored as success or failure.

\placeholder{AUTHOR CHECK --- measurement gap.} The termination categories do
not currently separate a model-initiated stop --- a \texttt{give\_up} call, or
a \texttt{done} call on an incomplete task --- from an episode that ended in
the wrong physical state after the agent kept acting. These are different
failure modes: the first is a termination-criterion failure, the second a
control failure, and distinguishing them is necessary for the failure analysis
in Section~\ref{sec:results}. We therefore record the terminal tool call and
its step index as a separate field.

Recovery attempts are archived separately, in fresh output directories, so that
the original attempt is preserved and a partially sent request is never
continued. Replacement outcomes follow the recorded recovery history, and the
analysis does not select the best outcome among retries. Changes to a financial
cap are documented because they change the protocol. Cost is reported twice:
for the retained outcomes, and summed across all attempts.

\subsection{Measurements and reproducibility}
Besides success we record logical model calls, simulator steps, input and
output tokens, cached input, reasoning tokens where available, API latency,
estimated cost, output truncations, and termination category, per call and
aggregated per episode. Transcripts and videos support inspection of object
identification, grasp acquisition and retention, collision or displacement,
placement, and premature declarations of completion. Explanatory notes are
treated as agent reports and checked against the observed trajectory rather
than taken as evidence about it.

Per-run metadata records configuration settings, task and state indices, seeds,
repository revisions, selected source hashes, package versions, prompts, and
cost assumptions. The representative setup uses \texttt{hf-libero 0.1.4},
\texttt{robosuite 1.4.0}, \texttt{MuJoCo 3.8.1}, and
\texttt{inspect-robots-agent 0.26.0}~\citep{inspectrobots}; the pinned
Inspect Robots revision is
\texttt{4d35fe4b81a643c0e61e8d287810b3c26f9d386a}. Saved requests, responses,
tool schemas, action trajectories, and videos provide an audit trail for both
the effective settings and the physical execution. The same records back the
interactive tables and rollout inspector on the project site.

% ===========================================================================
\section{Results}
\label{sec:results}
% ===========================================================================
The fifteen complete configurations comprise 1,500 episodes, of which 285
(19.0\%) satisfy LIBERO's success predicate. The campaign consumed 83{,}153
model calls, 915{,}160 scored simulator steps, and 190.0 hours of API latency,
at a settled cost of \$3{,}655. Because every configuration covers the same
ten tasks $\times$ ten initial states, the balanced rate of
Eq.~\eqref{eq:sr} equals the unweighted mean, and every comparison below is
matched on task and initial state.

\subsection{Success by model and reasoning effort}
\label{sec:res:models}
\placeholder{FIGURE: success rate by model and effort, with 95\% Wilson
intervals.}

\placeholder{Paragraph 1 --- effort is monotone for Astra
$0.28\to0.38\to0.44$, Sol $0.24\to0.30\to0.33$, and Opus~5.5
$0.09\to0.18\to0.33$. Opus~5 instead changes $0.10\to0.12\to0.06$.
Astra at high effort is the best configuration at $0.44$ [0.35, 0.54].}

\placeholder{Paragraph 2 --- one model fails completely: GPT-6 Luna records no
success in 300 episodes (95\% upper bound 0.04), and its behaviour differs by
effort --- 8.5 calls and 137 steps at low effort, against 67.5 and 74.0 calls
at medium and high, where it reaches the 80-call allowance in 23\% and 45\% of
episodes.}

\subsection{Success by task}
\label{sec:res:tasks}
\placeholder{FIGURE: success rate by task, pooled over the fifteen complete
configurations, 150 episodes per task.}

\placeholder{Paragraph --- the distribution is bimodal. Three basket tasks and
the stove task carry nearly all successes (0.29--0.52); three tasks are 0/150
across five models and three effort levels, which makes an environment-side
cause --- an unreachable pose, an occluded target, or a predicate the achieved
state does not satisfy --- as plausible as a policy-side one. Inspect the
rollouts for tasks 4, 6 and 9 and state which the videos support.}

\subsection{Inference cost}
\label{sec:res:cost}
\placeholder{TABLE or FIGURE: calls, tokens, latency and dollars per success,
per configuration.}

\placeholder{Paragraph --- accuracy and cost are close to independent. Sol and
Opus~5.5 both reach $0.33$ at high effort, at \$4.96 and \$5.46 per success;
Opus~5 at high effort costs \$94.57 per success. The spread in cost per success
across configurations that succeed at all is a factor of twenty, against a
factor of five in success rate.}

\subsection{How episodes end}
\label{sec:res:termination}
Of the 1,500 episodes, 285 (19.0\%) end in success, 1,115 (74.3\%) in an
unsuccessful terminal state, 93 (6.2\%) on a budget or limit stop, and 7
(0.5\%) on an API error. The limit stops are concentrated: 68 of the 93 belong
to GPT-6 Luna at medium and high effort, where the agent runs out of calls
without converging. Episodes that stop on a limit or an error are reported in
their own categories and counted as failures in every rate above; none is
silently dropped.

\placeholder{AUTHOR CHECK: the legacy \texttt{incomplete\_budget} label merges
the call limit and the dollar cap. Separate the 91 using the recorded call
count and settled cost --- ``the allowance bound the episode'' and ``spending
did'' have different implications.}

% ===========================================================================
\section{Conclusion}
% ===========================================================================
We have defined and justified a protocol for evaluating API-served multimodal
language models as closed-loop controllers on LIBERO-10: shared observations
and tools, environment-defined scoring rather than the agent's own declaration
of completion, a call allowance derived from the call usage of successful
trajectories and verified not to truncate successes, and per-episode inference
and trajectory records.

\placeholder{Close with the findings once Section~\ref{sec:results} is final:
the best zero-shot success rate, the model-dependent effect of requested effort, the
separation between accuracy and cost, the tasks no configuration solves, and
the model that never succeeds.}

Two limits on scope follow from the design. Because the interface is specific
--- two cameras, proprioception, general manipulation guidance, interpolated
absolute-pose tools, a 900-step horizon, and an explicit call allowance ---
success rates here are not comparable to published LIBERO scores obtained
under other observation or action interfaces, and we do not place them in the
same table. And because each (task, state) cell is observed once per
configuration, conclusions hold for the tested model versions under this
interface; extending them requires more tasks, more configurations, and
repeated trials per cell.

\section*{Acknowledgments}
\placeholder{Add acknowledgments for the final version.}

\clearpage
\appendix
% ===========================================================================
\section{Exact System Prompt}
\label{app:prompt}
% ===========================================================================
\IfFileExists{system_prompt.txt}{%
  \lstinputlisting{system_prompt.txt}%
}{%
  \placeholder{Place the archived prompt at system\_prompt.txt next to this
  file; the listing is omitted while that file is missing.}%
}

% ===========================================================================
\section{Task Inventory}
\label{app:tasks}
% ===========================================================================
Tasks are indexed 0--9 in the order listed.
\begin{enumerate}
\setcounter{enumi}{-1}
\item Put both the alphabet soup and the tomato sauce in the basket.
\item Put both the cream cheese box and the butter in the basket.
\item Turn on the stove and put the moka pot on it.
\item Put the black bowl in the bottom drawer of the cabinet and close it.
\item Put the white mug on the left plate and the yellow and white mug on the right plate.
\item Pick up the book and place it in the back compartment of the caddy.
\item Put the white mug on the plate and the chocolate pudding to the right of the plate.
\item Put both the alphabet soup and the cream cheese box in the basket.
\item Put both moka pots on the stove.
\item Put the yellow and white mug in the microwave and close it.
\end{enumerate}

\clearpage
\bibliography{experiment_references}

\end{document}
